\chapter{\texorpdfstring{Further values of $a$ with $\Om_n \cong \Ct{n}$: the results of Li}{Further
  values of a with the full wreath power: the results of Li}}
\label{chap:li}

This chapter formalizes the results of Li \cite{li2020, li2021} that extend
Theorem~\ref{thm:section3} to further values of $a$. Li's two papers use the notation of
Chapter~2 with the rescaled polynomial of Definition~\ref{def:normpoly}:
$g = |a| X^2 + \sgn(a)$, $\gamma_1 = 1$, $\gamma_{n+1} = g(\gamma_n)$, so that
$\gamma_n = |c_n|/|a|$ and $\beta_n = |b_n|$ for $n \ge 2$. (In \cite{li2020} the polynomial is
written $x^2 - a$ with $a > 0$; below, $a$ always denotes the coefficient of $f = X^2 + a$, so the
case of \cite{li2020} is $a < 0$, $a \equiv 3 \bmod 4$.) The Lean sources are in
\texttt{QuadraticIterates/Li/}.

\section{The results}

\begin{theorem}[\cite{li2021}, Theorem 3.3]
  \label{thm:li-A}
  \lean{QuadraticIterates.li2021_theorem_3_3, QuadraticIterates.nonempty_mulEquiv_of_forall_not_isSquare_abs_bSeq}
  \leanok
  \uses{def:galois, def:wreath}
  Let $a = -(8k+2)(8k+3)$ with $k \ge 0$. Then $\Om_n \cong \Ct{n}$ for all $n \ge 1$.
\end{theorem}
\begin{proof}
  \leanok
  \uses{lem:li-sqfree, lem:li-A-nonsqfree, thm:section1, def:normpoly}
  $-a = (8k+2)(8k+3) \equiv 6 \bmod 8$ is not a square. For $n \ge 2$, $|b_n| = \beta_n$ is not a
  square: by Lemma~\ref{lem:li-sqfree}~(2) if $n$ is squarefree, by
  Lemma~\ref{lem:li-A-nonsqfree} otherwise. Part~2) of Theorem~\ref{thm:section1}.
  \formalnote{The case split on \texttt{Squarefree n} in front of
  \lname{section1\_of\_not\_isSquare\_abs\_bSeq} is
  \lname{nonempty\_mulEquiv\_of\_forall\_not\_isSquare\_abs\_bSeq}, shared with
  Theorem~\ref{thm:li-B}.}
\end{proof}

\begin{theorem}[\cite{li2021}, Theorem 3.9]
  \label{thm:li-B}
  \lean{QuadraticIterates.li2021_theorem_3_9}
  \leanok
  \uses{def:galois, def:wreath}
  Let $a = -\bigl((4k+1)(4k+2) + 1\bigr)$ with $k \ge 0$. Then $\Om_n \cong \Ct{n}$ for all
  $n \ge 1$.
\end{theorem}
\begin{proof}
  \leanok
  \uses{lem:li-sqfree, lem:li-B-odd, lem:li-B-even, thm:section1, def:normpoly}
  $-a = (4k+1)(4k+2) + 1 \equiv 3 \bmod 4$ is not a square, and $a \equiv 1 \bmod 4$. For
  $n \ge 2$: if $n$ is squarefree, Lemma~\ref{lem:li-sqfree}~(3); if not and $4 \nmid n$,
  Lemma~\ref{lem:li-B-odd}; if $4 \dvd n$, then $n/\rad(n)$ is even and Lemma~\ref{lem:li-B-even}
  applies. Part~2) of Theorem~\ref{thm:section1}.
\end{proof}

\begin{theorem}[\cite{li2020}, Theorems 3.4 and 3.5]
  \label{thm:li-C}
  \lean{QuadraticIterates.li2020_theorem_3_5, QuadraticIterates.li2020_theorem_3_4, QuadraticIterates.bSeq_two}
  \leanok
  \uses{def:galois, def:wreath, def:bseq}
  Let $a < 0$ with $a \equiv 3 \bmod 4$ and $-a$ not a square. Then $|b_n|$ is not a square for
  all $n \ge 3$, and
  \[ \Om_n \cong \Ct{n} \text{ for all } n \ge 1 \iff -a - 1 \text{ is not a square in } \ZZ. \]
\end{theorem}
\begin{proof}
  \leanok
  \uses{lem:li-C-nonsqfree, lem:li-C-sqfree, thm:section1, def:normpoly, def:twoindep}
  For $n \ge 3$, $|b_n| = \beta_n$ is not a square: Lemma~\ref{lem:li-C-nonsqfree} if $n$ is not
  squarefree, Lemma~\ref{lem:li-C-sqfree} if it is (then $n > 2$). Moreover
  $b_2 = c_2/b_1 = -(a+1) = -a - 1 > 0$. If $-a - 1$ is not a square, no $|b_n|$ with $n \ge 2$
  is, and part~2) of Theorem~\ref{thm:section1} gives $\Om_n \cong \Ct{n}$ for all $n$. If
  $-a - 1 = b_2$ is a square, then $b_1, b_2$ are not $2$-independent, so by part~1) of
  Theorem~\ref{thm:section1} (the equivalence (a) $\Leftrightarrow$ (c) at $n = 2$),
  $\Om_2 \not\cong \Ct{2}$.
  \formalnote{The value $b_2 = -a - 1$ is \lname{bSeq\_two}, from $c_2 = b_1 b_2$
  (\lname{cSeq\_eq\_prod\_bSeq}). The direction ``$\Rightarrow$'' uses \lname{section1\_tfae}
  and \lname{TwoIndependent.not\_isSquare}. \cite{li2020} also shows that in the exceptional case
  $-a - 1 = m^2$ the index of $\Om_n$ in $\Ct{n}$ is exactly $2$ for all $n \ge 2$; this needs an
  analysis of the Kummer tower when the base is not maximal, which the present infrastructure
  (Lemma~\ref{lem:1.5} assumes $\Om_n \cong \Ct{n}$) does not provide, and it is not formalized.}
\end{proof}

The first two theorems cover $a = -6, -110, -342, \dots$ and $a = -3, -31, -91, \dots$; the third
covers $a = -13, -21, -29, \dots$ and identifies the failures $a = -5, -17, -37, \dots$ (where
$b_2 = -a - 1$ is a square). All three are proved, as Theorem~\ref{thm:section3} is, by showing
that no $|b_n|$ with $n \ge 2$ (resp.\ $n \ge 3$) is a square and invoking part~2) of
Theorem~\ref{thm:section1}; what changes is how the non-squareness of $\beta_n$ is established.
The three theorems are also part of the comparator setup of the repository, one challenge
statement each.

\emph{Numerical check.} All congruences and identities below were verified in PARI/GP for the
first members of each family and $n \le 13$ before the formalization, and the lists of $n$ with
$|b_n|$ a square were computed: empty for $a = -6, -110, -342, -3, -31, -91, -13, -21, -29, -33$
and $\{2\}$ for $a = -5, -17, -37$. The check exposed a misprint in \cite{li2021}, Lemma~2.2: its
third case must read $a \equiv 1 \bmod 4$, not $3 \bmod 4$ (the case $a \equiv 3 \bmod 4$ is the
one of \cite{li2020}, and $a = -5$ is a counterexample to the printed statement).

\section{The criterion of Lemma 2.1, one index at a time}

The proof of Lemma~\ref{lem:2.1} has three layers, which the formalization keeps apart because
Li's arguments reuse the two lower ones with different inputs: the arithmetic of the Möbius
product ($\beta_n$ as a quotient $P/Q$ of two products of the $\gamma_{kt}$, and the criterion
$P \equiv -Q \bmod m$), the congruence patterns that produce $P \equiv -Q$ (values of
$\gamma_{kt}$ depending only on a class of $t$, with a sign partition balanced on each class),
and the sequence facts feeding the patterns (eventual constancy on the multiples of $k$ modulo
$\gamma_k + \gamma_{2k}$; $2$-periodicity modulo $\gamma_k + \gamma_{k+2}$). The criteria of
\cite{li2021}, Section~2, are then four short lemmas (\ref{lem:li-crit}--\ref{lem:li-odd}), each
stated for a single index $n$ with its own modulus, and Lemmas~\ref{lem:2.1} and~\ref{lem:2.2}
are corollaries.

Throughout, $g$ is an even integer polynomial, $\varepsilon = \pm 1$, all $\gamma_n$ ($n \ge 1$)
are nonzero, and for $n \ge 2$ we write $n' := \rad(n)$ and $k := n/n'$; ``$m$ is a modulus''
means $m \ge 1$, and ``$x$ is a unit mod $m$'' means that $x$ and $m$ are coprime. Li always
takes $m$ to be a prime $p \equiv 3 \bmod 4$; the formalization keeps the modulus form of
Lemma~\ref{lem:2.1}, which needs no prime and turns ``$p \nmid \gamma$'' into coprimality. The
sign partition of the divisors $t$ of $n'$ is the pair of sets
$S_\pm := \{t \dvd n' : \mu(n'/t) = \pm 1\}$ (Lemma~\ref{lem:moebiussign}).

\begin{lemma}[Balanced classes]
  \label{lem:li-classes}
  \lean{Finset.prod_eq_neg_prod_of_forall_card_filter_eq, Finset.card_filter_eq_one_eq_card_filter_eq_neg_one_of_sum_eq_zero}
  \leanok
  Let $R$ be a commutative ring, $S_+$ and $S_-$ disjoint finite sets with $a \in S_+ \cup S_-$,
  $x \colon S_+ \cup S_- \to R$, $w \colon I \to R$ and $\kappa \colon S_+ \cup S_- \to I$ a
  ``class'' function such that $x_t = w_{\kappa(t)}$ for $t \ne a$ and $x_a = -w_{\kappa(a)}$. If
  every class is balanced, i.e.\ $\#\{t \in S_+ : \kappa(t) = i\} = \#\{t \in S_- : \kappa(t) = i\}$
  for all $i$, then
  \[ \prod_{t \in S_+} x_t = -\prod_{t \in S_-} x_t. \]
\end{lemma}
\begin{proof}
  \leanok
  Group each product by classes (\lname{Finset.prod\_comp}): on the class $i$ the product over
  $S_\pm$ is $w_i^{\#\{t \in S_\pm : \kappa(t) = i\}}$, times $-1$ if $a \in S_\pm$. The exponents
  agree by balance, and exactly one of $S_+$, $S_-$ contains $a$.
  \formalnote{With $I$ a one-point set this gives \lname{prod\_gammaSeq\_mul\_eq\_neg\_prod}
  (Lemma~\ref{lem:li-crit}), with $I = \ZZ/2$ the parity classes of Lemma~\ref{lem:li-odd}. The
  companion \lname{Finset.card\_filter\_eq\_one\_eq\_card\_filter\_eq\_neg\_one\_of\_sum\_eq\_zero}
  produces the balance: a $\pm 1$-valued function summing to $0$ on a finite set takes each sign
  equally often.}
\end{proof}

\begin{lemma}[Parity-balanced sign partition]
  \label{lem:li-parity}
  \lean{ArithmeticFunction.sum_divisors_filter_mod_two_moebius_div, ArithmeticFunction.sum_divisors_filter_dvd_moebius_div, Squarefree.card_filter_moebius_div_eq_one_eq_card_filter_eq_neg_one}
  \leanok
  \uses{lem:moebiussign}
  Let $n \ge 3$. Then
  \[ \sum_{t \dvd n,\ t \text{ odd}} \mu(n/t) = 0
     = \sum_{t \dvd n,\ t \text{ even}} \mu(n/t). \]
  Consequently, for squarefree $n' \ge 3$, the sign partition $S_+ \sqcup S_-$ of the divisors of
  $n'$ is balanced on the odd and on the even divisors separately.
\end{lemma}
\begin{proof}
  \leanok
  The sum over the even divisors is the restricted Möbius sum $\sum_{t \dvd n,\ 2 \dvd t} \mu(n/t)$,
  which is $[n = 2] = 0$ (Lemma~\ref{lem:moebiussign},
  \lname{sum\_divisors\_filter\_dvd\_moebius\_div}); the sum over the odd divisors is the
  difference of the total sum $\sum_{t \dvd n} \mu(n/t) = 0$ and the even one. For squarefree
  $n'$ the values $\mu(n'/t)$ are $\pm 1$, so on any set of divisors with vanishing sum the two
  signs occur equally often (\lname{Squarefree.card\_filter\_moebius\_div\_eq\_one\_eq\_card\_filter\_eq\_neg\_one}).
  \formalnote{The hypothesis $n \ge 3$ only excludes $n = 2$, where the odd class $\{1\}$ is
  unbalanced; no squarefreeness is needed for the vanishing of the sums. The residue classes
  are written as $t \bmod 2 = i$ with $i$ ranging over $\NN$ (the classes $i \ge 2$ are empty),
  which is the form the class function $\kappa(t) = t \bmod 2$ of Lemma~\ref{lem:li-classes}
  consumes.}
\end{proof}

\begin{lemma}[$\beta_n$ from a numerator/denominator congruence]
  \label{lem:li-PQ}
  \lean{QuadraticIterates.not_isSquare_betaSeq_of_prod_eq_neg_prod}
  \leanok
  \uses{def:gammaseq, lem:moebiussign, lem:zmodsquare}
  Let $n \ge 1$, $n' := \rad(n)$, $k := n/n'$, and put $P := \prod_{t \in S_+} \gamma_{kt}$,
  $Q := \prod_{t \in S_-} \gamma_{kt}$. Suppose $P = c P'$ and $Q = c Q'$ with an integer
  $c \ne 0$, and let $m$ be a modulus with $-1$ not a square mod $m$ such that
  $P' \equiv -Q' \bmod m$ and $Q'$ is a unit mod $m$. Then $\beta_n$ is not a square in $\QQ$.
\end{lemma}
\begin{proof}
  \leanok
  $\beta_n = \prod_{d \dvd n} \gamma_d^{\mu(n/d)} = P/Q = P'/Q'$ by
  Lemma~\ref{lem:moebiussign} (\lname{prod\_pow\_moebius\_eq\_div}), and a quotient of two
  integers with $P' \equiv -Q'$, $Q'$ a unit, is not a rational square when $-1$ is not a square
  mod $m$ (Lemma~\ref{lem:zmodsquare}, \lname{ZMod.isSquare\_neg\_one\_of\_isSquare\_div}).
  \formalnote{The common factor $c$ is $1$ everywhere except in Lemma~\ref{lem:li-C-sqfree},
  where it is a power of $\gamma_2$ (Li works modulo a higher power of the modulus instead). The
  hypothesis ``$P' \equiv -Q'$ with $Q'$ a unit'' is stated in $\ZZ/m\ZZ$.}
\end{proof}

\begin{lemma}[\cite{li2021}, Lemma 2.1: Stoll's criterion at one index]
  \label{lem:li-crit}
  \lean{QuadraticIterates.not_isSquare_betaSeq_of_dvd_add_two_mul, QuadraticIterates.prod_gammaSeq_mul_eq_neg_prod, QuadraticIterates.isUnit_prod_gammaSeq_mul}
  \leanok
  \uses{def:gammaseq}
  Let $n \ge 2$, $k := n/\rad(n)$, and let $m$ be a modulus with $m \dvd \gamma_k + \gamma_{2k}$,
  $m$ coprime to $\gamma_k$, and $-1$ not a square mod $m$. Then $\beta_n$ is not a square in
  $\QQ$.
\end{lemma}
\begin{proof}
  \leanok
  \uses{lem:li-classes, lem:li-PQ, lem:moebiussign}
  Exactly the proof of Lemma~\ref{lem:2.1}: mod $m$, $\gamma_{tk} \equiv \gamma_{2k}$ for all
  $t \ge 2$ (\lname{gammaSeq\_mul\_eq\_two\_mul}) and $\gamma_k \equiv -\gamma_{2k}$, so
  Lemma~\ref{lem:li-classes} with a single class gives $P \equiv -Q$
  (\lname{prod\_gammaSeq\_mul\_eq\_neg\_prod}), $Q$ is a unit because
  $\gamma_{2k} \equiv -\gamma_k$ is (\lname{isUnit\_prod\_gammaSeq\_mul}), and
  Lemma~\ref{lem:li-PQ} concludes.
  \formalnote{Lemma~\ref{lem:2.1} (\lname{not\_isSquare\_betaSeq}) is recovered by instantiating
  the modulus at the index $k = n/\rad(n)$.}
\end{proof}

\begin{lemma}[\cite{li2021}, Corollary 2.4]
  \label{lem:li-succ}
  \lean{QuadraticIterates.not_isSquare_betaSeq_of_dvd_add_succ, QuadraticIterates.not_isSquare_abs_bSeq_of_dvd_add_succ}
  \leanok
  \uses{def:gammaseq}
  Let $g(0)$ be a unit, $n \ge 2$, $k := n/\rad(n)$, and $m$ a modulus with
  $m \dvd \gamma_k + \gamma_{k+1}$ and $-1$ not a square mod $m$. Then $\beta_n$ is not a square
  in $\QQ$.
\end{lemma}
\begin{proof}
  \leanok
  \uses{lem:li-crit}
  $\gamma_k + \gamma_{k+1}$ divides $\gamma_k + \gamma_{2k}$ (\lname{gammaSeq\_add\_succ\_dvd})
  and is coprime to $\gamma_k$ (\lname{isCoprime\_gammaSeq\_add\_succ}); both properties pass to
  the divisor $m$, and Lemma~\ref{lem:li-crit} applies.
  \formalnote{Lemma~\ref{lem:2.2} goes through this statement
  (\lname{not\_isSquare\_betaSeq\_of\_pos\_of\_not\_isSquare\_neg\_one} is the case
  $m = \gamma_k + \gamma_{k+1}$). The per-index form is what Theorem~\ref{thm:li-C} needs: there
  $\gamma_1 + \gamma_2 = -a \equiv 1 \bmod 4$ fails the hypothesis, but every
  $\gamma_k + \gamma_{k+1}$ with $k \ge 2$ satisfies it.}
\end{proof}

\begin{lemma}[$2$-periodicity]
  \label{lem:li-period}
  \lean{QuadraticIterates.gammaSeq_eq_ite_even_of_add_two_eq_zero, QuadraticIterates.gammaSeq_mul_eq_ite_even_of_odd, QuadraticIterates.gammaSeq_eq_ite_even_of_add_two_eq}
  \leanok
  \uses{def:gammaseq}
  Over any commutative ring, if $k \ge 1$ and $\gamma_k + \gamma_{k+2} = 0$, then for all
  $r \ge k + 1$
  \[ \gamma_r = \begin{cases} -\gamma_k & \text{if } r \equiv k \bmod 2, \\
                              \gamma_{k+1} & \text{otherwise.} \end{cases} \]
  In particular, for odd $k$ and $t \ge 2$, $\gamma_{kt} = -\gamma_k$ if $t$ is odd and
  $\gamma_{kt} = \gamma_{k+1}$ if $t$ is even.
\end{lemma}
\begin{proof}
  \leanok
  $g(-\gamma_k) = g(\gamma_k) = \gamma_{k+1}$ since $g$ is even, and
  $g(\gamma_{k+1}) = \gamma_{k+2} = -\gamma_k$; so the sequence alternates between the two values
  from index $k + 1$ on. For odd $k$, $kt \equiv k \bmod 2$ exactly when $t$ is odd.
  \formalnote{The same induction, without evenness, gives the $2$-periodicity from a $2$-cycle
  $\gamma_{n_0+2} = \gamma_{n_0}$ (\lname{gammaSeq\_eq\_ite\_even\_of\_add\_two\_eq}), used for
  the residue computations of Section~\ref{sec:li-residues}. All are applied in $\ZZ/m\ZZ$
  through \lname{intCast\_gammaSeq}, as the existing congruence lemmas are.}
\end{proof}

\begin{proposition}[\cite{li2021}, Proposition 2.5, $k$ even]
  \label{lem:li-even}
  \lean{QuadraticIterates.not_isSquare_betaSeq_of_even_of_dvd_add_two, QuadraticIterates.not_isSquare_abs_bSeq_of_even_of_dvd_add_two, isCoprime_of_isCoprime_gcd_of_dvd_add, IsRelPrime.of_gcd_right_of_dvd_add}
  \leanok
  \uses{def:gammaseq}
  Let $n \ge 2$ with $k = n/\rad(n)$ even, and $m$ a modulus with $m \dvd \gamma_k + \gamma_{k+2}$,
  $m$ coprime to $\gamma_2$, and $-1$ not a square mod $m$. Then $\beta_n$ is not a square.
\end{proposition}
\begin{proof}
  \leanok
  \uses{lem:li-period, lem:li-crit, lem:strongdvd}
  Mod $m$, Lemma~\ref{lem:li-period} gives $\gamma_{2k} \equiv -\gamma_k$ ($2k \equiv k \bmod 2$),
  so $m \dvd \gamma_k + \gamma_{2k}$. Strong divisibility (\lname{gammaSeq\_associated\_gcd},
  $\gcd(k, k+2) = 2$) gives $\gcd(\gamma_k, \gamma_{k+2}) = \gamma_2$; a common divisor of $m$
  and $\gamma_k$ divides $\gamma_{k+2}$, hence $\gamma_2$, hence is $1$
  (\lname{isCoprime\_of\_isCoprime\_gcd\_of\_dvd\_add}, over any Bézout domain).
  Lemma~\ref{lem:li-crit}.
\end{proof}

\begin{proposition}[\cite{li2021}, Proposition 2.5, $k$ odd]
  \label{lem:li-odd}
  \lean{QuadraticIterates.not_isSquare_betaSeq_of_odd_of_dvd_add_two, QuadraticIterates.not_isSquare_abs_bSeq_of_odd_of_dvd_add_two, QuadraticIterates.prod_gammaSeq_mul_eq_neg_prod_of_odd, QuadraticIterates.isUnit_prod_gammaSeq_mul_of_odd, QuadraticIterates.isCoprime_gammaSeq_of_odd_of_dvd_add_two}
  \leanok
  \uses{def:gammaseq}
  Let $g(0)$ be a unit, $n \ge 2$ with $k = n/\rad(n)$ odd and $k > 1$, and $m$ a modulus with
  $m \dvd \gamma_k + \gamma_{k+2}$, $m$ coprime to $\gamma_{k+1}$, and $-1$ not a square mod $m$.
  Then $\beta_n$ is not a square.
\end{proposition}
\begin{proof}
  \leanok
  \uses{lem:li-period, lem:li-parity, lem:li-classes, lem:li-PQ, lem:strongdvd}
  Since $k$ is odd and $> 1$, $n' = \rad(n)$ is neither $1$ nor $2$ (an odd prime factor of $k$
  divides $n'$). Mod $m$, by Lemma~\ref{lem:li-period}, $\gamma_{kt} \equiv -\gamma_k$ for odd
  $t \ge 3$ and $\gamma_{kt} \equiv \gamma_{k+1}$ for even $t$. Both values are units mod $m$:
  $\gamma_{k+1}$ by hypothesis, $\gamma_k$ because $\gcd(\gamma_k, \gamma_{k+2}) = \gamma_1$ is a
  unit ($\gcd(k, k+2) = 1$), so a common divisor of $m$ and $\gamma_k$ divides $\gamma_{k+2}$ and
  hence $\gamma_1$ (\lname{isCoprime\_gammaSeq\_of\_odd\_of\_dvd\_add\_two}). Take the
  parity-balanced partition of Lemma~\ref{lem:li-parity} and apply Lemma~\ref{lem:li-classes}
  with the two classes ``odd'' ($w = -\gamma_k$; at $t = 1$ the value $\gamma_k = -w$) and
  ``even'' ($w = \gamma_{k+1}$): $P \equiv -Q$ with $Q$ a unit
  (\lname{prod\_gammaSeq\_mul\_eq\_neg\_prod\_of\_odd},
  \lname{isUnit\_prod\_gammaSeq\_mul\_of\_odd}). Lemma~\ref{lem:li-PQ}.
  \formalnote{Li derives the coprimality of $p$ and $\gamma_{k+1}$ from $p$ being an odd prime;
  in modulus form it is a hypothesis, discharged for the rescaled polynomial by
  Lemma~\ref{lem:li-cop}.}
\end{proof}

\section{Residues of the rescaled sequence}
\label{sec:li-residues}

In this section $a < 0$, $g = |a| X^2 - 1$ is the rescaled polynomial of
Definition~\ref{def:normpoly}, so $\gamma_1 = 1$, $\gamma_2 = |a| - 1$, and all $\gamma_n$ are
positive. The residue computations are done in $\ZZ/8\ZZ$ or $\ZZ/4\ZZ$ by the $2$-periodicity
of Lemma~\ref{lem:li-period}.

\begin{lemma}[The sequence modulo $\gamma_2^2$]
  \label{lem:li-modsq}
  \lean{QuadraticIterates.gammaSeq_two_sq_dvd_sub_ite_even, QuadraticIterates.gammaSeq_two_dvd_of_even}
  \leanok
  \uses{def:normpoly}
  For every even $g$ with $g(0)^2 = 1$ and $\gamma_1 = 1$, and all $n \ge 2$:
  $\gamma_n \equiv \gamma_2$ for even $n$ and $\gamma_n \equiv g(0) \bmod \gamma_2^2$ for odd $n$.
\end{lemma}
\begin{proof}
  \leanok
  By induction. If $\gamma_n \equiv \gamma_2$ then $\gamma_2 \dvd \gamma_n$, so
  $\gamma_{n+1} = g(\gamma_n) \equiv g(0)$ because $\gamma_n^2 \dvd g(\gamma_n) - g(0)$
  (\lname{sq\_dvd\_gammaSeq\_succ\_sub}); if $\gamma_n \equiv g(0)$ then
  $\gamma_n^2 \equiv g(0)^2 = 1$, so $\gamma_{n+1} = g(\gamma_n) \equiv g(1) = \gamma_2$
  (\lname{EvenPoly.dvd\_eval\_sub}).
  \formalnote{Used mod $\gamma_2$ in Lemma~\ref{lem:li-C-sqfree} and mod $(4k+1)^2$ in
  Lemma~\ref{lem:li-B-even}.}
\end{proof}

\begin{lemma}[$a$ even]
  \label{lem:li-mod8-2}
  \lean{QuadraticIterates.gammaSeq_normPoly_zmod_eight_of_even}
  \leanok
  \uses{def:normpoly}
  If $a < 0$ is even, then $\gamma_n \equiv -a - 1 \bmod 8$ for all $n \ge 2$; in particular
  $\gamma_n \equiv 5 \bmod 8$ when $a \equiv 2 \bmod 8$.
\end{lemma}
\begin{proof}
  \leanok
  In $\ZZ/8\ZZ$, $\gamma_2 = g(1) = -a - 1$ is odd, so $\gamma_2^2 = 1$ and $g(\gamma_2) = g(1)$:
  a fixed point (\lname{gammaSeq\_eq\_eval\_one}).
\end{proof}

\begin{lemma}[$a \equiv 1 \bmod 4$]
  \label{lem:li-mod8-1}
  \lean{QuadraticIterates.gammaSeq_normPoly_zmod_eight_of_emod_four_eq_one, QuadraticIterates.gammaSeq_normPoly_add_two_emod_eight_eq_six, QuadraticIterates.gammaSeq_normPoly_emod_four_eq_two_of_even}
  \leanok
  \uses{def:normpoly}
  If $a < 0$ and $a \equiv 1 \bmod 4$, then $\gamma_n \equiv 3 \bmod 8$ for odd $n \ge 3$ and
  $\gamma_n \equiv -a - 1 \bmod 8$ (which is $2$ or $6$) for even $n \ge 2$. Consequently
  $\gamma_k + \gamma_{k+2} \equiv 6 \bmod 8$ for odd $k \ge 3$, and $\gamma_k \equiv 2 \bmod 4$ for
  even $k \ge 2$.
\end{lemma}
\begin{proof}
  \leanok
  \uses{lem:li-period}
  In $\ZZ/8\ZZ$, $|a| \in \{3, 7\}$, $\gamma_2 = |a| - 1 \in \{2, 6\}$, $\gamma_3 = |a| \cdot 4 - 1
  = 3$, $\gamma_4 = |a| \cdot 9 - 1 = |a| - 1 = \gamma_2$, and the pattern repeats with period $2$
  (Lemma~\ref{lem:li-period}, $2$-cycle form). The two evaluations are \texttt{decide} after a
  case split on $|a| \bmod 8$.
\end{proof}

\begin{lemma}[$a \equiv 3 \bmod 4$]
  \label{lem:li-mod4-3}
  \lean{QuadraticIterates.gammaSeq_normPoly_zmod_four_of_emod_four_eq_three, QuadraticIterates.gammaSeq_normPoly_add_succ_emod_four_eq_three, QuadraticIterates.gammaSeq_eq_ite_even_of_eval_one_eq_zero}
  \leanok
  \uses{def:normpoly}
  If $a < 0$ and $a \equiv 3 \bmod 4$, then $\gamma_n \equiv 0 \bmod 4$ for even $n \ge 2$ and
  $\gamma_n \equiv 3 \bmod 4$ for odd $n \ge 3$; hence $\gamma_k + \gamma_{k+1} \equiv 3 \bmod 4$
  for all $k \ge 2$.
\end{lemma}
\begin{proof}
  \leanok
  \uses{lem:li-period}
  In $\ZZ/4\ZZ$, $g = X^2 - 1$, $\gamma_2 = 0$, $\gamma_3 = -1$, $\gamma_4 = 0$, and the pattern
  repeats.
  \formalnote{The pattern $1, 0, g(0), 0, g(0), \dots$ holds over any ring for an even $g$ with
  $g(1) = 0$, $g(0)^2 = 1$ and $\gamma_1 = 1$ (\lname{gammaSeq\_eq\_ite\_even\_of\_eval\_one\_eq\_zero});
  the same lemma is the core of Lemma~\ref{lem:li-cop}.}
\end{proof}

\begin{lemma}[Coprimality at odd indices]
  \label{lem:li-cop}
  \lean{QuadraticIterates.dvd_two_of_dvd_add_two_of_dvd_succ, QuadraticIterates.isCoprime_gammaSeq_succ_of_odd_of_add_two_eq_two_mul}
  \leanok
  \uses{def:normpoly}
  Let $g$ be even with $g(0)^2 = 1$ and $\gamma_1 = 1$, and let $k \ge 3$ be odd. Then every
  common divisor of $\gamma_k + \gamma_{k+2}$ and $\gamma_{k+1}$ divides $2$; in particular, if
  $\gamma_k + \gamma_{k+2} = 2M$ with $M$ odd, then $M$ is coprime to $\gamma_{k+1}$.
\end{lemma}
\begin{proof}
  \leanok
  \uses{lem:li-mod4-3}
  Let $d$ be a common divisor and compute in $\ZZ/d\ZZ$: $\gamma_{k+1} = 0$, so
  $\gamma_{k+2} = g(0)$ and $\gamma_k = -g(0)$; then
  $0 = \gamma_{k+1} = g(\gamma_k) = g(-g(0)) = g(g(0)) = g(1)$ (evenness and $g(0)^2 = 1$). For
  such a $g$ the sequence is $1, 0, g(0), 0, g(0), \dots$ (Lemma~\ref{lem:li-mod4-3}, general
  form), so $\gamma_k = g(0)$ for the odd $k \ge 3$; together with $\gamma_k = -g(0)$ this gives
  $2 = 0$ in $\ZZ/d\ZZ$.
  \formalnote{Applied to $g = |a| X^2 - 1$ with $a \equiv 1 \bmod 4$ in Lemma~\ref{lem:li-B-odd},
  where $\gamma_k + \gamma_{k+2} \equiv 6 \bmod 8$ makes $M \equiv 3 \bmod 4$.}
\end{proof}

\section{The applications}

\begin{lemma}[\cite{li2021}, Lemma 2.2: squarefree indices]
  \label{lem:li-sqfree}
  \lean{QuadraticIterates.not_isSquare_abs_bSeq_of_squarefree, QuadraticIterates.not_isSquare_abs_bSeq_of_squarefree_of_pos_of_emod_eight_eq_four, QuadraticIterates.not_isSquare_abs_bSeq_of_squarefree_of_neg_of_emod_eight_eq_two, QuadraticIterates.not_isSquare_abs_bSeq_of_squarefree_of_neg_of_emod_four_eq_one}
  \leanok
  \uses{def:normpoly}
  Let $n > 1$ be squarefree. Then $|b_n|$ is not a square as soon as $-1$ is not a square modulo
  $\gamma_1 + \gamma_2 = |a| + \sgn(a) + 1$; this holds in each of the cases
  (1) $a > 0$, $a \equiv 4 \bmod 8$; (2) $a < 0$, $a \equiv 2 \bmod 8$; (3) $a < 0$,
  $a \equiv 1 \bmod 4$.
\end{lemma}
\begin{proof}
  \leanok
  \uses{lem:li-succ}
  Here $k = 1$, so Lemma~\ref{lem:li-succ} with $m := \gamma_1 + \gamma_2 = 1 + g(1)$ applies as
  soon as $-1$ is not a square mod $m$: $m = a + 2 \equiv 6 \bmod 8$ in case (1), $m = -a \equiv
  6 \bmod 8$ in case (2), and $m = -a \equiv 3 \bmod 4$ in case (3)
  (\lname{ZMod.not\_isSquare\_neg\_one\_of\_emod\_eight\_eq\_six},
  \lname{\dots\_of\_emod\_four\_eq\_three}).
  \formalnote{Case (3) is the misprinted one, see above. Stoll's cases $a > 0$,
  $a \equiv 1, 2 \bmod 4$ and $a < 0$, $a \equiv 0 \bmod 4$ fit the same statement; they are not
  repeated here because Lemma~\ref{lem:2.2} already gives them for all $n$.}
\end{proof}

\subsection*{Theorem A}

\begin{lemma}[\cite{li2021}, Lemma 3.1]
  \label{lem:li-A-nonsqfree}
  \lean{QuadraticIterates.not_isSquare_abs_bSeq_of_not_squarefree_of_eq_neg_mul}
  \leanok
  \uses{def:normpoly}
  Let $a = -(8k+2)(8k+3)$ and let $n \ge 1$ be not squarefree. Then $|b_n|$ is not a square.
\end{lemma}
\begin{proof}
  \leanok
  \uses{lem:li-succ, lem:li-mod8-2}
  Let $k_0 := n/\rad(n) \ge 2$. Since $|a|\,y^2 + y - 1 = \bigl((8k+3)y - 1\bigr)\bigl((8k+2)y +
  1\bigr)$, the number $m := (8k+2)\gamma_{k_0} + 1$ divides
  $\gamma_{k_0} + \gamma_{k_0+1} = |a|\gamma_{k_0}^2 + \gamma_{k_0} - 1$, and
  $m \equiv 2 \cdot 5 + 1 \equiv 3 \bmod 4$ by Lemma~\ref{lem:li-mod8-2}. Lemma~\ref{lem:li-succ}.
  \formalnote{Here $a \equiv 2 \bmod 8$ because $(8k+2)(8k+3) \equiv 6 \bmod 8$. The
  discriminant $1 - 4a = (16k+5)^2$ is what makes the quadratic in $y$ split; \cite{li2021},
  Remark~3.2, notes that the complementary family $-(8k+5)(8k+6)$ is out of reach of this
  method ($a = -30$: no prime $\equiv 3 \bmod 4$ divides $\gamma_2 + \gamma_3$).}
\end{proof}

\subsection*{Theorem B}

\begin{lemma}[\cite{li2021}, Lemma 3.4]
  \label{lem:li-B-odd}
  \lean{QuadraticIterates.not_isSquare_abs_bSeq_of_not_squarefree_of_not_four_dvd_of_emod_four_eq_one, Nat.even_div_radical_iff_four_dvd, Nat.two_le_div_radical_of_not_squarefree, Nat.squarefree_iff_radical_eq_self}
  \leanok
  \uses{def:normpoly}
  Let $a < 0$, $a \equiv 1 \bmod 4$, and let $n \ge 1$ be not squarefree with $4 \nmid n$. Then
  $|b_n|$ is not a square.
\end{lemma}
\begin{proof}
  \leanok
  \uses{lem:li-odd, lem:li-mod8-1, lem:li-cop}
  $k = n/\rad(n)$ is odd (as $4 \nmid n$; \lname{Nat.even\_div\_radical\_iff\_four\_dvd}) and
  $> 1$ (as $n$ is not squarefree), so $k \ge 3$. By Lemma~\ref{lem:li-mod8-1},
  $\gamma_k + \gamma_{k+2} \equiv 6 \bmod 8$, so $M := (\gamma_k + \gamma_{k+2})/2 \equiv 3 \bmod 4$;
  it is coprime to $\gamma_{k+1}$ by Lemma~\ref{lem:li-cop}. Proposition~\ref{lem:li-odd}.
\end{proof}

\begin{lemma}[\cite{li2021}, Lemma 3.6]
  \label{lem:li-B-even}
  \lean{QuadraticIterates.not_isSquare_abs_bSeq_of_even_div_radical_of_eq_neg_mul_add_one}
  \leanok
  \uses{def:normpoly}
  Let $a = -\bigl((4k+1)(4k+2) + 1\bigr)$, $A := 4k+1$, $B := 4k+2$, $\alpha := -a = AB + 1$, and
  let $n \ge 1$ be such that $k_0 := n/\rad(n)$ is even. Then $|b_n|$ is not a square.
\end{lemma}
\begin{proof}
  \leanok
  \uses{lem:li-even, lem:li-mod8-1, lem:li-modsq, lem:strongdvd}
  \emph{The factorization.} In $\ZZ[y]$, with $g(y) = \alpha y^2 - 1$,
  \[ y + g(g(y))
     = (\alpha y + A)\,\bigl(\alpha^2 y^3 - A\alpha\, y^2 - (\alpha + B)\, y + B\bigr) \]
  (a polynomial identity in $k$ and $y$; the quotient was found by polynomial division in
  PARI). Hence $\alpha\gamma_{k_0} + A \dvd \gamma_{k_0} + \gamma_{k_0+2}$.

  \emph{The modulus.} $\gamma_2 = AB$ divides $\gamma_{k_0}$ (strong divisibility, $2 \dvd k_0$),
  and $\alpha \equiv 1 \bmod A$, so $A \dvd \alpha\gamma_{k_0} + A$; put
  $m := (\alpha\gamma_{k_0} + A)/A$. Then $m \equiv 3 \bmod 4$: $\gamma_{k_0} \equiv 2 \bmod 4$
  (Lemma~\ref{lem:li-mod8-1}), $\alpha \equiv 3 \bmod 4$, $A \equiv 1 \bmod 4$.

  \emph{Coprimality with $\gamma_2 = AB$.} Modulo $A^2$: $\gamma_{k_0} \equiv \gamma_2 = AB
  \equiv A \bmod A^2$ (Lemma~\ref{lem:li-modsq} reduced from $\gamma_2^2$ to $A^2$, and
  $AB = A^2 + A$), so $\alpha\gamma_{k_0} + A \equiv (\alpha + 1)A$ and $m \equiv \alpha + 1 \equiv
  2 \bmod A$, coprime to the odd $A$. Modulo $B = 2(2k+1)$: $m$ is odd, and since
  $B \dvd \gamma_{k_0}$, $A(m - 1) = \alpha\gamma_{k_0} \equiv 0 \bmod (2k+1)$ with
  $A \equiv -1 \bmod (2k+1)$, so $m \equiv 1 \bmod (2k+1)$.

  Proposition~\ref{lem:li-even} with this $m$.
  \formalnote{The discriminant $4\alpha - 3 = (8k+3)^2$ is what makes $y + g(g(y))$ split off a
  linear factor; \cite{li2021}, Remark~3.8, records that $-((4k+2)(4k+3)+1)$ (e.g.\ $a = -507$)
  is out of reach. The lemma covers $k_0 = 2$ as well, so \cite{li2021}, Lemma~3.5 (the case
  $k_0 = 2$ for arbitrary $a \equiv 1 \bmod 4$, with the modulus $(\gamma_2 + \gamma_4)/(2\gamma_2)$)
  is not needed and is not formalized; note that its printed claim
  $\gcd((\gamma_2+\gamma_4)/\gamma_2, \gamma_2) = 1$ is false (the gcd is $2$, e.g.\ $a = -3$:
  $182$ and $2$), and the odd part is what the argument uses. The Lean proof follows the three
  paragraphs above: the factorization is \lname{add\_eval\_eval\_eq\_mul}, the modulus with
  its residue is \lname{exists\_mul\_eq\_and\_emod\_four\_eq\_three\_of\_even}, its
  divisibility property \lname{dvd\_add\_two\_of\_mul\_eq}, and the coprimality argument is
  the Bézout-style \lname{isCoprime\_of\_mul\_eq}, applied in
  \lname{isCoprime\_gammaSeq\_two\_of\_mul\_eq}.}
\end{proof}

\subsection*{Theorem C}

\begin{lemma}[\cite{li2020}, Lemma 3.1]
  \label{lem:li-C-nonsqfree}
  \lean{QuadraticIterates.not_isSquare_abs_bSeq_of_not_squarefree_of_neg_of_emod_four_eq_three}
  \leanok
  \uses{def:normpoly}
  Let $a < 0$, $a \equiv 3 \bmod 4$, and let $n \ge 1$ be not squarefree. Then $|b_n|$ is not a
  square.
\end{lemma}
\begin{proof}
  \leanok
  \uses{lem:li-succ, lem:li-mod4-3}
  $k = n/\rad(n) \ge 2$, so $\gamma_k + \gamma_{k+1} \equiv 3 \bmod 4$ (Lemma~\ref{lem:li-mod4-3});
  Lemma~\ref{lem:li-succ} with $m := \gamma_k + \gamma_{k+1}$.
\end{proof}

\begin{lemma}[\cite{li2020}, Lemmas 3.2 and 3.3]
  \label{lem:li-C-sqfree}
  \lean{QuadraticIterates.not_isSquare_abs_bSeq_of_squarefree_of_neg_of_emod_four_eq_three, QuadraticIterates.not_isSquare_betaSeq_of_squarefree_of_not_isSquare_neg_one}
  \leanok
  \uses{def:normpoly}
  Let $a < 0$, $a \equiv 3 \bmod 4$, and let $n > 2$ be squarefree. Then $|b_n|$ is not a square.
\end{lemma}
\begin{proof}
  \leanok
  \uses{lem:li-classes, lem:li-parity, lem:li-PQ, lem:li-modsq}
  Here $k = 1$; take $m := \gamma_2 = |a| - 1$, which is $\equiv 0 \bmod 4$, so $-1$ is not a
  square mod $m$ (\lname{ZMod.not\_isSquare\_neg\_one\_of\_four\_dvd}). By
  Lemma~\ref{lem:li-modsq}, every even divisor $t$ of $n$ has $\gamma_t = m\,u_t$ with
  $u_t \equiv 1 \bmod m$, and every odd divisor $t \ge 3$ has $\gamma_t \equiv -1 \bmod m$;
  $\gamma_1 = 1$. With the parity-balanced partition of Lemma~\ref{lem:li-parity} ($n > 2$),
  $P = m^N P'$ and $Q = m^N Q'$ with $N$ the common number of even divisors in $S_+$ and $S_-$,
  and $P'$, $Q'$ the products of the $\gamma_t$ ($t$ odd) and $u_t$ ($t$ even).
  Lemma~\ref{lem:li-classes} with the classes ``odd'' ($w = -1$) and ``even'' ($w = 1$) gives
  $P' \equiv -Q'$ with $Q' \equiv \pm 1$; Lemma~\ref{lem:li-PQ} with $c = m^N$.
  \formalnote{Li treats odd $n$ (all divisors odd; his Lemma~3.2, modulo $4$) and even $n$
  (Lemma~3.3, modulo a power of $\gamma_2$) separately; modulo $\gamma_2$ the argument is the
  same for both, which is why the two lemmas are one declaration. The general statement
  (\lname{not\_isSquare\_betaSeq\_of\_squarefree\_of\_not\_isSquare\_neg\_one}) holds for every
  even $g$ with $g(0) = -1$ and $\gamma_1 = 1$, with the modulus $\gamma_2$.}
\end{proof}
